% GENERATED FILE. Edit canonical lessons, then run npm run generate:books.
% canonical-source-hash: fnv1a64:70a0f31a73d23ca5
% canonical-lessons: TE-C203-tittu, TE-C203-potladu, TE-C203-edugu, TE-C203-puyu, TE-C203-kuriyu

\chapter{To scold, To fight, To grow up, To bloom, To rain down}
\label{ch:te-to-scold-to-fight-to-grow-up-to-bloom-to-rain-down}
\hlchaptermodality{203}

\begin{chapteropening}
\textbf{By the end of this chapter:} \emph{I can say to scold, to fight, to grow up, to bloom and to rain down.}

Complete the last lesson of chapter 203: I can say to scold, to fight, to grow up, to bloom and to rain down.
\end{chapteropening}

\section[\texorpdfstring{\te{తిట్టు} (tiṭṭu)}{తిట్టు (tiṭṭu)}]{\te{తిట్టు} (tiṭṭu) — to scold, to curse}

\label{lesson:TE-C203-tittu}



\begin{quote}
Before the new one: say the Telugu for to register, then the Telugu for to practise.
\end{quote}

\subsection*{You'll want to know: \te{తిట్టు}}
\textbf{\te{తిట్టు}} — \emph{tiṭṭu} — \textquotedblleft{}to scold, to curse\textquotedblright{}.

\begin{grammarlens}[title={What you've built}]
One more word for this chapter.
\end{grammarlens}

\subsection*{Your turn}
\begin{itemize}
\raggedright
  \item \emph{Say it:} \emph{tiṭṭu}
  \item \emph{Say it:} \emph{tiṭṭu}, once more
  \item \emph{Say it:} say \emph{tiṭṭu}
  \item \emph{Recall:} say the Telugu for to register, then the Telugu for to practise, then say \emph{tiṭṭu} again
\end{itemize}

\begin{tcolorbox}[breakable,colback=teal!4,colframe=teal!35!black,title={Before you move on}]
What does \te{తిట్టు} mean? (\textquotedblleft{}To scold, to curse\textquotedblright{}.) Say it once more.
\end{tcolorbox}
\section[\texorpdfstring{\te{పోట్లాడు} (pōṭlāḍu)}{పోట్లాడు (pōṭlāḍu)}]{\te{పోట్లాడు} (pōṭlāḍu) — to fight, to quarrel}

\label{lesson:TE-C203-potladu}



\begin{quote}
Before the new one: say the Telugu for to practise, then the Telugu for to scold.
\end{quote}

\subsection*{You'll want to know: \te{పోట్లాడు}}
\textbf{\te{పోట్లాడు}} — \emph{pōṭlāḍu} — \textquotedblleft{}to fight, to quarrel\textquotedblright{}.

\begin{grammarlens}[title={What you've built}]
One more word for this chapter.
\end{grammarlens}

\subsection*{Your turn}
\begin{itemize}
\raggedright
  \item \emph{Say it:} \emph{pōṭlāḍu}
  \item \emph{Say it:} \emph{pōṭlāḍu}, once more
  \item \emph{Say it:} say \emph{pōṭlāḍu}
  \item \emph{Recall:} say the Telugu for to practise, then the Telugu for to scold, then say \emph{pōṭlāḍu} again
\end{itemize}

\begin{tcolorbox}[breakable,colback=teal!4,colframe=teal!35!black,title={Before you move on}]
What does \te{పోట్లాడు} mean? (\textquotedblleft{}To fight, to quarrel\textquotedblright{}.) Say it once more.
\end{tcolorbox}
\section[\texorpdfstring{\te{ఎదుగు} (edugu)}{ఎదుగు (edugu)}]{\te{ఎదుగు} (edugu) — to grow up}

\label{lesson:TE-C203-edugu}



\begin{quote}
Before the new one: say the Telugu for to scold, then the Telugu for to fight.
\end{quote}

\subsection*{You'll want to know: \te{ఎదుగు}}
\textbf{\te{ఎదుగు}} — \emph{edugu} — \textquotedblleft{}to grow up\textquotedblright{}.

\begin{grammarlens}[title={What you've built}]
One more word for this chapter.
\end{grammarlens}

\subsection*{Your turn}
\begin{itemize}
\raggedright
  \item \emph{Say it:} \emph{edugu}
  \item \emph{Say it:} \emph{edugu}, once more
  \item \emph{Say it:} say \emph{edugu}
  \item \emph{Recall:} say the Telugu for to scold, then the Telugu for to fight, then say \emph{edugu} again
\end{itemize}

\begin{tcolorbox}[breakable,colback=teal!4,colframe=teal!35!black,title={Before you move on}]
What does \te{ఎదుగు} mean? (\textquotedblleft{}To grow up\textquotedblright{}.) Say it once more.
\end{tcolorbox}
\section[\texorpdfstring{\te{పూయు} (pūyu)}{పూయు (pūyu)}]{\te{పూయు} (pūyu) — to bloom, to flower}

\label{lesson:TE-C203-puyu}



\begin{quote}
Before the new one: say the Telugu for to fight, then the Telugu for to grow up.
\end{quote}

\subsection*{You'll want to know: \te{పూయు}}
\textbf{\te{పూయు}} — \emph{pūyu} — \textquotedblleft{}to bloom, to flower\textquotedblright{}.

\begin{grammarlens}[title={What you've built}]
One more word for this chapter.
\end{grammarlens}

\subsection*{Your turn}
\begin{itemize}
\raggedright
  \item \emph{Say it:} \emph{pūyu}
  \item \emph{Say it:} \emph{pūyu}, once more
  \item \emph{Say it:} say \emph{pūyu}
  \item \emph{Recall:} say the Telugu for to fight, then the Telugu for to grow up, then say \emph{pūyu} again
\end{itemize}

\begin{tcolorbox}[breakable,colback=teal!4,colframe=teal!35!black,title={Before you move on}]
What does \te{పూయు} mean? (\textquotedblleft{}To bloom, to flower\textquotedblright{}.) Say it once more.
\end{tcolorbox}
\section[\texorpdfstring{\te{కురియు} (kuriyu)}{కురియు (kuriyu)}]{\te{కురియు} (kuriyu) — to rain down, to pour (of rain)}

\label{lesson:TE-C203-kuriyu}



\begin{quote}
Before the new one: say the Telugu for to grow up, then the Telugu for to bloom.
\end{quote}

\subsection*{You'll want to know: \te{కురియు}}
\textbf{\te{కురియు}} — \emph{kuriyu} — \textquotedblleft{}to rain down, to pour (of rain)\textquotedblright{}.

\begin{grammarlens}[title={What you've built}]
One more word for this chapter.
\end{grammarlens}

\subsection*{Your turn}
\begin{itemize}
\raggedright
  \item \emph{Say it:} \emph{kuriyu}
  \item \emph{Say it:} \emph{kuriyu}, once more
  \item \emph{Say it:} say \emph{kuriyu}
  \item \emph{Recall:} say the Telugu for to grow up, then the Telugu for to bloom, then say \emph{kuriyu} again
\end{itemize}

\begin{tcolorbox}[breakable,colback=teal!4,colframe=teal!35!black,title={Before you move on}]
What does \te{కురియు} mean? (\textquotedblleft{}To rain down, to pour (of rain)\textquotedblright{}.) Say it once more.
\end{tcolorbox}
